% !PRINTED 161
\weberpage{161}

Wird also
\weq{3}{\dfrac{\sigma'(u)}{\sigma(u)} = \zeta(u)}
gesetzt, so ist $\zeta(u)$ eine eindeutige Funktion von $u$, und zwar
ein elliptisches Integral zweiter Gattung:
\weq{4}{\zeta(u) = -\int \wp(u)\,du = -\int \dfrac{\wp\,d\wp}{\sqrt{4\wp^3 - g_2\wp - g_3}} \ [\S\,46,\,(19)],}
worin die additive Konstante dadurch bestimmt ist, daß $\zeta(u)$
eine \gesperrt{ungerade} Funktion sein muß. Die Periodizität der $\zeta$-Funktion
 ergibt sich aus den Periodengleichungen der $\sigma$-Funktion
[\S\,35, (9)]:
\weq{5}{\begin{aligned} \zeta(u+\omega_1) - \zeta(u) &= 2\eta_1, \\ \zeta(u+\omega_2) - \zeta(u) &= 2\eta_2. \end{aligned}}
Es sind also $2\eta_1$ und $2\eta_2$ als vollständige Integrale zweiter
Gattung (analog den $E$, $E'$) dargestellt:
\weq{6}{2\eta_1 = -\int\limits_u^{u+\omega_1} \wp(u)\,du, \quad 2\eta_2 = -\int\limits_u^{u+\omega_2} \wp(u)\,du}
% [sic] The source duplicates the article in “mit der der Legendre[schen] Relation”.
mit der der \gesperrt{Legendre}schen Relation entsprechenden Gleichung
\weq{7}{\eta_1\omega_2 - \eta_2\omega_1 = \pi i \ [\S\,35,\,(10)].}
Um die Additionstheoreme für die \gesperrt{Weierstrass}schen
Transzendenten zu bilden, gehen wir von der $\sigma$-Funktion aus.
Für diese erhalten wir aus (1) mit Benutzung der Additionsformel
 für $\vartheta_{11}$, \S\,22, (4):
\weq{8}{\dfrac{\sigma(u+v)\sigma(u-v)}{\sigma(u)^2\sigma(v)^2} = \wp(v) - \wp(u),}
daraus durch logarithmische Differentiation:
\weq{9}{\begin{aligned} \zeta(u+v) + \zeta(u-v) - 2\zeta(u) &= \dfrac{\wp'(u)}{\wp(u) - \wp(v)}, \\ \zeta(u+v) - \zeta(u-v) - 2\zeta(v) &= \dfrac{-\wp'(v)}{\wp(u) - \wp(v)}, \end{aligned}}
\weq{10}{\zeta(u+v) = \zeta(u) + \zeta(v) + \dfrac{1}{2}\dfrac{\wp'(u) - \wp'(v)}{\wp(u) - \wp(v)}.}
Hieraus leitet man durch abermalige Differentiation das
Additionstheorem für die $\wp$-Funktion her. Man erhält aus (9)
durch Differentiation nach $u$ und $v$:
%CONT>>
